% Multinotes test file for exercises

% Toggle the following options to see their effect
%\PassOptionsToPackage{noanswers}{multinotes}
%\PassOptionsToPackage{blankanswers}{multinotes}

%------------------------------------------------
\documentclass{article}
%------------------------------------------------

% paper
\usepackage[margin=30mm, a4paper]{geometry}
\setlength{\parindent}{0ex}
\setlength{\parskip}{1ex}

% multinotes
\usepackage{../multinotes}
\answercolour{purple}
\framedanswers
\setboxrule{0.5pt}

% misc
\usepackage{lipsum}
\usepackage{fancyvrb}

% meta
\title{Typesetting exercises with {\tt multinotes.sty}}
\author{DE}

%------------------------------------------------
\begin{document}
\maketitle

%----------------------------------------
\section{Answers to exercises}
%----------------------------------------

The \verb+\ans+ macro and \verb+answer+ environment respond to the \verb+blankanswers+ and \verb+noanswers+ options.

\begin{itemize}
\item \verb+blankanswers+ replaces the content by blank space. 
\item \verb+noanswers+ removes the answer completely.
\end{itemize}

Here is an answer environment (removed with \verb+noanswers+ option).
\begin{answer}
\lipsum[1]
\end{answer}

This is produced by the following code.
\begin{Verbatim}[frame=single]
\begin{answer}
\lipsum[1]
\end{answer}
\end{Verbatim}

Here is an answer environment with a custom title (removed with \verb+noanswers+ option).

\begin{answer}[Solution]
\lipsum[1]
\end{answer}

This is produced by the following code.
\begin{Verbatim}[frame=single]
\begin{answer}[Solution]
\lipsum[1]
\end{answer}
\end{Verbatim}

%----------------------------------------
\section{Questions, parts and subparts}
%------------------------------------------------

The package replicates list environments found in {\tt exam.cls} (namely {\tt questions}, {\tt parts}, {\tt subparts}, {\tt choices} and {\tt checkboxes}), and makes them responsive to the {\tt blankanswers} and {\tt noanswers} options. For example,

\begin{questions}
\question First question.
\begin{parts}
\part First part. \ans{Answer to first part.}
\part Second part.
\begin{subparts}
\subpart First subpart. \ans{Answer to first subpart.}
\subpart Second subpart. \ans{Answer to second subpart.}
\end{subparts}
\part Third part. \ans{Answer to third part.}
\end{parts}
\question Second question. \ans{Answer to second question.}
\end{questions}

This is produced by the following code.
\begin{Verbatim}[frame=single]
\begin{questions}
\question First question.
\begin{parts}
\part First part. \ans{Answer to first part.}
\part Second part.
\begin{subparts}
\subpart First subpart. \ans{Answer to first subpart.}
\subpart Second subpart. \ans{Answer to second subpart.}
\end{subparts}
\part Third part. \ans{Answer to third part.}
\end{parts}
\question Second question. \ans{Answer to second question.}
\end{questions}
\end{Verbatim}


%
%--------------------
\section{Multiple choice questions}

\begin{questions}
\question
In what year did Columbus first cross the Atlantic?
\begin{choices}
\choice 1490 
\choice 1491 
\correctchoice 1492 
\choice 1493 
\end{choices}
\question
The rain in Spain falls mainly on the plain.
\begin{choices}
\correctchoice True 		
\choice False		
\end{choices}
\end{questions}

This is produced by the following code.
\begin{Verbatim}[frame=single]
\begin{questions}
\question
In what year did Columbus first cross the Atlantic?
\begin{choices}
\choice 1490 
\choice 1491 
\correctchoice 1492 
\choice 1493 
\end{choices}
\question
The rain in Spain falls mainly on the plain.
\begin{choices}
\correctchoice True 		
\choice False		
\end{choices}
\end{questions}
\end{Verbatim}

%--------------------
\section{Multiple answer questions}

Which of the following were members of the Beatles?
\begin{checkboxes}
\correctchoice John
\correctchoice Paul
\correctchoice George
\choice Zippy
\end{checkboxes}

This is produced by the following code.
\begin{Verbatim}[frame=single]
Which of the following were members of the Beatles?
\begin{checkboxes}
\correctchoice John
\correctchoice Paul
\correctchoice George
\choice Zippy
\end{checkboxes}
\end{Verbatim}

%--------------------
\section{Fill the blanks}

\begin{questions}
\question Roses are \fillbox{red}, violets are \fillbox{blue}.
\question State Einstein's equation: $E=$\fillbox{$mc^2$}.
\end{questions}

This is produced by the following code.
\begin{Verbatim}[frame=single]
\begin{questions}
\question Roses are \fillbox{red}, violets are \fillbox{blue}.
\question State Einstein's equation: $E=$\fillbox{$mc^2$}.
\end{questions}
\end{Verbatim}


%------------------------------------------------
\section{Example: prime numbers quiz}
%------------------------------------------------

\begin{questions}

\question
What are the prime factors of 42?
\begin{choices}
\choice $3\times 14$
\correctchoice $2\times 3\times 7$
\choice $2\times 21$
\end{choices}

\question
What are the prime factors of 150?
\begin{choices}
\correctchoice $2\times 3\times 5\times 5$
\choice $3\times 3\times 5\times 5$
\choice $6\times 25$
\end{choices}

\question
What is the LCM of 13 and 3?
\begin{choices}
\choice $13$
\choice $3$
\correctchoice $39$
\end{choices}

\question
What is the HCF of 32 and 24?
\begin{choices}
\choice $2$
\correctchoice $8$
\choice $4$
\end{choices}

\question
What is the LCM of 24 and 36?
\begin{choices}
\choice $12$
\correctchoice $72$
\choice $30$
\end{choices}

\question
What is the HCF of 104 and 136?
\begin{choices}
\correctchoice $8$
\choice $12$
\choice $13$
\end{choices}

\question
What is the HCF of $3\times 13$ and $2^2\times 13$?
\begin{choices}
\choice $9$
\choice $17$
\correctchoice $13$
\end{choices}

\question
What is the LCM of $2^2\times 3^2\times 5$ and $2^3\times 3\times 5$?
\begin{choices}
\correctchoice $360$
\choice $16\,120$
\choice $30$
\end{choices}

\question
What is the HCF of $2\times 5\times 13$ and $22\times 13$?
\begin{choices}
\choice $13$
\correctchoice $26$
\choice $52$
\end{choices}

\question
Find the LCM of $3^2\times 11\times 13$ and $3\times 11^2$.
\begin{choices}
\choice $1287$
\correctchoice $14\,157$
\choice $429$
\end{choices}

\end{questions}


%------------------------------------------------
\end{document}
%------------------------------------------------


